% Generated draft: requires mathematical and source review.
\documentclass[11pt]{article}
\usepackage[margin=25mm]{geometry}
\usepackage{fontspec}
\setmainfont{Latin Modern Roman}
\usepackage{amsmath,amssymb,mathtools,unicode-math}
\setmathfont{Latin Modern Math}
\usepackage{enumitem,needspace,xurl,hyperref}
\setlength{\parindent}{0pt}
\setlength{\parskip}{5pt}
\setlength{\emergencystretch}{3em}
\newcommand{\N}{\mathbb{N}}
\newcommand{\Z}{\mathbb{Z}}
\newcommand{\Q}{\mathbb{Q}}
\newcommand{\R}{\mathbb{R}}
\newcommand{\C}{\mathbb{C}}
\newcommand{\F}{\mathbb{F}}
\newcommand{\A}{\mathbb{A}}
\newcommand{\PP}{\mathbb{P}}
\newcommand{\E}{\mathbb{E}}
\newcommand{\eps}{\varepsilon}
\newcommand{\dd}{\,\mathrm{d}}
\begin{document}
\section*{UnsolvedMath ID: 20000423}
\textbf{Source problem code:} catalogue entry
\subsection{Definitions and mathematical statement}
W. Raskind's source question asks for the isomorphism
\[
 \operatorname{End}_K(A)\otimes_{\Z}\Q_p\longrightarrow
 \operatorname{End}_{G_K}(V_p(A))
\]
when $K/\Q_p$ is finite and $A/K$ has multiplicative reduction. Here $G_K=\operatorname{Gal}(\overline K/K)$, $T_p(A)=\varprojlim A[p^n](\overline K)$ and $V_p(A)=T_p(A)\otimes\Q_p$; the target is the algebra of $\Q_p$-linear equivariant endomorphisms. Endomorphisms on the left must be defined over $K$ to give this map. In higher dimension the report treats totally degenerate toric reduction, split in its examples; the original does not specify that convention.

The supplied report refutes general surjectivity. It attributes to Gregory--Liedtke, for every prime $p\geq5$ with $p\equiv1\pmod3$, an abelian surface $B/\Q_p$ with split totally degenerate reduction for which
\[
 \dim_{\Q_p}(\operatorname{End}_{\Q_p}(B)\otimes\Q_p)=1,\qquad
 \dim_{\Q_p}\operatorname{End}_{G_{\Q_p}}(V_p(B))=2.
\]
Thus rationalizing does not remove the failure. The report isolates a nonscalar equivariant idempotent in the target not induced by an algebraic endomorphism. Consequently it claims $B$ is $\Q_p$-isogeny-simple but $V_p(B)$ is decomposable as a Galois representation.

In its period notation the projector is written
\[
 M=\begin{pmatrix}a^2&ab\\ab&b^2\end{pmatrix}.
\]
The record does not define $a,b$, their normalization or the realization of this $2\times2$ matrix on the four-dimensional $V_p(B)$. For a nonzero matrix of this form, idempotence requires $a^2+b^2=1$; that missing normalization and the representation coordinates must be supplied before this becomes an explicit verified endomorphism. The existence claim is retained without guessing these data.

For multiplicative elliptic curves, the report claims the original map is an isomorphism by an extension-class argument. Hence dimension two is the first possible positive-dimensional failure within this reduction class, according to the report.

The original comments suggest a global analogue for varieties having multiplicative reduction somewhere, suggest Serre's elliptic methods, and attribute a Drinfel'd-modular case to Ito. These are historical source comments with no full statements or citations. A local counterexample does not imply failure of an analogous global endomorphism theorem.
\subsection{Short English statement}
Ask whether all Galois-equivariant endomorphisms of a $p$-adic Tate module come from algebraic endomorphisms for multiplicative reduction. The report gives abelian-surface counterexamples with a two-dimensional centralizer but only scalar algebraic endomorphisms, while retaining the elliptic case. Its advertised explicit projector lacks the parameters and coordinates needed for verification.
\subsection{Sources}
\begin{itemize}
\item Local registry: \url{data/unsolvedmath/problems.json}, permanent ID $20000423$, source code \texttt{catalogue entry}; original question, all comments, supplied numbering correction and research summary.
\item ``Problems from the AIM Tate Conjecture Workshop'', July 23--27, 2007, transcribed by Christopher Lyons, W. Raskind's item 12 according to supplied verification (extracted as 2): \url{https://aimath.org/WWN/tateconjecture/tateconjecture.pdf}. Canonical location \texttt{aim-algebraic-number-theory-notes.json}, \texttt{notes[74]}. No PDF was opened here.
\item Original comments name Serre and Ito and a Drinfel'd-modular case. The summary attributes counterexamples to Gregory--Liedtke; no full bibliographic references or period data are supplied. Original literature field empty.
\item Supplied machine-generated research attempt 1: counterexample; status refuted at the run, registry solved and classification \texttt{SOLVED-IN-LITERATURE}; rigor metadata ``proved'', novelty confidence low. The non-peer-reviewed projector synthesis requires independent verification before citation.
\end{itemize}
\end{document}
